\section{Give every state word an index and an address}
Arrange the sixteen state words as four rows of four words. Number indices from
zero in row order. The address of word \(i\) is \(B+4i\), where \(B\) is the state
base. A row occupies sixteen bytes. A column is a set of four words sixteen bytes
apart; it is not four adjacent words.

\begin{table}[htbp]
\centering\caption{State fields in row order. Byte offsets are decimal.}
\tagpdfsetup{table/header-rows={1}}
\begin{tabular}{rlll}\toprule
Indices & Offsets & Contents & Size\\\midrule
0--3 & 0--15 & Four fixed constant words & 16 bytes\\
4--11 & 16--47 & Key bytes decoded as words & 32 bytes\\
12 & 48--51 & Initial block counter & 4 bytes\\
13--15 & 52--63 & Nonce bytes decoded as words & 12 bytes\\
\bottomrule\end{tabular}
\end{table}

The constants are \texttt{61707865}, \texttt{3320646e}, \texttt{79622d32} and
\texttt{6b206574}. In little-endian memory their bytes spell the fixed 16-byte
string \texttt{expand 32-byte k}. This string is a prescribed constant, not a
password supplied by the user. The key and nonce follow in their byte order.

For key bytes \texttt{00 01 02 03}, the first decoded key word is
\texttt{03020100}. On this little-endian target a word-aligned \texttt{lw} reads that
pattern directly. Compare this with the preceding SHA project: that algorithm's
source words required a different byte interpretation. Reusing its byte-combining
sequence unchanged here would reverse every input word.

Keep the state used to start a block unchanged until feed-forward is complete.
The scalar version copies it into the output area and uses that area as working
storage. The vector version retains an original copy in vector registers. Both
arrangements leave the caller's input available for later blocks and for inspection.
